\documentclass{ximera}


\title{Antiderivative Practice Problems}   % Use xmlatex all -s to compile when SERVE refuses to work.
\author{Kelly Stady}
\license{CC: 0}         % replace with an appropriate license, or set it in xmPreamble

\begin{document}
\begin{abstract}
    Time to practice!
\end{abstract}
\maketitle
\label{xim:AntiderivEx}



\begin{problem}
Evaluate the indefinite integral using the basic integration rules.


\[
\int (x^4 + e^{7x} - 3) \ dx = \frac{ \answer{x^5} }{ \answer{5} } + \frac{ e^{\answer{7x}} }{\answer{7}} - \answer{3x}  + C
\]

\end{problem}



\begin{problem}
Evaluate the indefinite integral by first simplifying the integrand, then use the power rule.

\[
\int \frac{x^3 - 5x + 6}{x^2} \ dx
\]

Rewrite the integrand in the appropriate form to find the antiderivative. 

\[
\frac{x^3 - 5x + 6}{x^2} = \answer{x} - \frac{\answer{5}}{\answer{x}} + 6 x^{\answer{-2}}
\]

Integrate term by term to get the final answer. 
\[
\int \frac{x^3 - 5x + 6}{x^2} \, dx = \frac{\answer{x^2}}{\answer{2}} - 5 \answer{\ln|x|} - \frac{\answer{6}}{\answer{x}} + C
\]
\end{problem}


\begin{problem}
Evaluate the indefinite integral using a substitution. 


\[
\int x^2 e^{x^3} \ dx
\]
\begin{enumerate}
\item Identify $u$ and find the derivative in differential form.

Let $u = \answer{x^3}$

Then $du = \answer{3x^2} \ dx$

\item Identify the correct integral rewritten in terms of $u.$

\begin{multipleChoice}
\choice{$\displaystyle \int 3 e^u \ du$}
\choice{$\displaystyle \int u^2 e^u \ du$}
\choice[correct]{$\displaystyle \int \frac{1}{3} e^u \ du$}
\choice{$\displaystyle \int e^u \ du$}
\end{multipleChoice}


\item Evaluate.

\[
\int x^2 e^{x^3} \, dx = \frac{ e^{\answer{x^3}} }{ \answer{3} } + C
\]

\end{enumerate}
\end{problem}

\begin{problem}
Evaluate the definite integral using a substitution and change the limits of integration.
\[
\int_{0}^{1} x(x^2 + 1)^3 \, dx
\]

Let $u = x^2 + 1$. 
Then $du = \answer{2x} \, dx$.

When $x=0$, the lower limit becomes $u = \answer{1}$.
When $x=1$, the upper limit becomes $u = \answer{2}$.

The evaluation yields:
\[
\int_{0}^{1} x(x^2 + 1)^3 \, dx = \answer{\frac{15}{8}}
\]
\end{problem}

\begin{problem}
Evaluate the indefinite integral that results in an inverse trigonometric function:
\[
\int \frac{1}{1 + 9x^2} \, dx
\]

Let $u = \answer{3x}$. 
Then $du = \answer{3} \, dx$.

The final antiderivative is:
\[
\int \frac{1}{1 + 9x^2} \, dx = \answer{\frac{1}{3}\arctan(3x)} + C
\]
\end{problem}

\begin{problem}
Follow the steps below to evaluate the indefinite integral.

\[
\int \frac{\sin x + \cos^3 x }{\cos^2 x} \, dx
\]

\begin{enumerate}
\item First, split the numerator of the integrand into two separate terms,

\[
\int \frac{\sin x}{\cos^2 x} + \answer{\cos x}  \ dx
\]

\item The first term can be rewritten as follows.  Simplify this term using standard trigonometric identities.

\[ \frac{\sin x}{\cos x} \cdot \frac{1}{\cos x} = \answer{\tan x} \cdot \answer{\sec x} \]

\item Lastly, rewrite the simplified integral and integrate term-by-term.
\[
\int \answer{\sec x \tan x} + \answer{\cos x} \ dx = \answer{\sec x} + \answer{\sin x} + C
\]

\end{enumerate}
\end{problem}


\begin{problem}
Follow the steps below to evaluate the definite integral.

 \[ \int_{0}^{\pi/4} \sec^2 x \ dx \]

\begin{enumerate}
\item State the antiderivative.  Do not include $C.$  

    \[ F(x) = \answer{tan x} \]


\item Find the value of the definite integral.

    \[ \int_{0}^{\pi/4} \sec^2 x \ dx = \answer{1}   \]

\end{enumerate}
\end{problem}


\begin{problem}
A model rocket is launched vertically upward from the ground ($s(0) = 0$ meters) at time $t = 0$ seconds.  
Its velocity function is given by 

\[
v(t) = 30t^2 + 4t
\]

where $v(t)$ is measured in meters per second.

\begin{enumerate}
\item Find the general antiderivative to determine the position function $s(t).$

\[
s(t) = \answer{10t^3 + 2t^2} + C
\]

\item Use the initial condition $s(0) = 0$ to solve for the exact position function.

\[
s(t) = \answer{10t^3 + 2t^2}
\]

\item How far has the rocket traveled after exactly $2$ seconds?

\[
s(2) = \answer{88} \text{ meters}
\]

\end{enumerate}
\end{problem}



\end{document}


